\chapter{Application Development}
\label{chap:ch3}

\par The application development phase encompasses the design and implementation of a comprehensive Digital Twin system for intelligent irrigation management. This chapter details the architectural decisions, technological stack, and implementation strategies employed to create a real-time simulation and control platform capable of comparing traditional threshold-based irrigation with data-driven machine learning approaches. The system integrates soil dynamics simulation, sensor data processing, and intelligent decision-making algorithms to optimize water usage in agricultural scenarios.


\section{Problem Statement and Background}
\label{sec:ch3sec1}

\subsection{Irrigation Management Challenges}

Irrigation management in agriculture faces critical challenges in balancing water conservation with crop yield optimization. Traditional threshold-based controllers use fixed soil moisture levels to trigger irrigation events, often resulting in either under-utilization of water resources or excessive consumption. Key issues include:

\begin{itemize}
    \item \textbf{Static decision thresholds}: Fixed moisture thresholds cannot adapt to dynamic environmental conditions (temperature, humidity, rainfall patterns).
    \item \textbf{Sensor sparsity}: Physical sensor networks are expensive to deploy and maintain, limiting real-time decision capability in resource-constrained environments.
    \item \textbf{Prediction uncertainty}: Traditional systems cannot anticipate weather events or account for complex soil-plant-atmosphere interactions.
    \item \textbf{Water waste}: Conventional irrigation often applies uniform water schedules regardless of actual crop water demand.
\end{itemize}

\subsection{Digital Twin Approach}

A Digital Twin provides a virtual replica of the physical irrigation system, enabling continuous monitoring, predictive analysis, and optimization. This thesis implements a Digital Twin that simulates soil moisture dynamics, evaluates irrigation strategies, and compares:

\begin{enumerate}
    \item \textbf{Baseline Threshold Control}: A standard hysteresis-based controller using moisture thresholds.
    \item \textbf{Sparse Sampling Optimization}: A reduced-sensor approach that minimizes measurement frequency while maintaining decision accuracy.
    \item \textbf{Intelligent ANFIS-GA Control}: An adaptive neuro-fuzzy inference system trained via genetic algorithms to learn optimal irrigation policies from historical data.
\end{enumerate}


\section{Research Design and Analysis}
\label{sec:ch3sec2}

\subsection{System Architecture Overview}

The Digital Twin system follows a modular three-tier architecture:

\begin{itemize}
    \item \textbf{Backend Simulation Layer (Python)}: Core simulation engines for soil dynamics, Digital Twin state management, and experiment execution.
    \item \textbf{Machine Learning Layer (Python)}: ANFIS model training, genetic algorithm optimization, and predictive inference.
    \item \textbf{Frontend Visualization Layer (SAPUI5)}: Real-time dashboard for parameter configuration, experiment execution monitoring, and results visualization.
\end{itemize}

\subsection{Experimental Features}

The system implements three distinct experimental modes to evaluate irrigation control strategies:

\subsubsection{Baseline Threshold Experiment}

A traditional proportional-integral control strategy using soil moisture thresholds with hysteresis:

\begin{equation}
\text{irrigation\_active} = \begin{cases}
\text{True} & \text{if } m(t) < \theta_{\text{lower}} \text{ and } r(t) = 0 \\
\text{False} & \text{if } m(t) > \theta_{\text{upper}} \text{ or } r(t) = 1 \\
\text{previous state} & \text{otherwise}
\end{cases}
\end{equation}

where $m(t)$ is soil moisture, $\theta_{\text{lower}} = \theta - \frac{h}{2}$ and $\theta_{\text{upper}} = \theta + \frac{h}{2}$ are hysteresis bounds, $\theta$ is the moisture threshold, $h$ is hysteresis width, and $r(t)$ indicates rain prediction.

\subsubsection{Sparse Sampling Experiment}

Compares full sensor measurements against reduced-frequency sampling to quantify information loss:

\begin{equation}
\text{measurement frequency} = \frac{1}{\text{sample\_interval}}
\end{equation}

The experiment tracks decision mismatches between continuous and sparse monitoring to establish the minimum sampling interval that maintains control accuracy.

\subsubsection{ANFIS-GA Intelligent Control}

Employs an Adaptive Neuro-Fuzzy Inference System trained to learn optimal irrigation policies:

\begin{equation}
\hat{p}(t) = \text{ANFIS}(\text{moisture}(t), \text{temperature}(t), \text{humidity}(t), \text{rain\_prediction}(t))
\end{equation}

where $\hat{p}(t) \in [0,1]$ is the predicted irrigation probability, and irrigation activates if $\hat{p}(t) \geq 0.1$.

\subsection{System Models}

\subsubsection{Soil Moisture Dynamics}

The soil state is modeled as a continuous moisture level (0-100\%):

\begin{equation}
m(t+1) = \text{clamp}\left(m(t) - \text{evap}(t) + r_{\text{amt}}(t) + w_{\text{irr}}(t), 0, 100\right)
\end{equation}

where:
\begin{itemize}
    \item Evapotranspiration: $\text{evap}(t) = 0.35 + 1.2 \cdot \frac{T(t)}{36} \cdot \frac{100 - H(t)}{100}$ (temperature and humidity dependent)
    \item Rain contribution: $r_{\text{amt}}(t) = U(1.5, 6.0)$ liters if rain occurs
    \item Irrigation contribution: $w_{\text{irr}}(t) = f_{\text{ml/min}} \cdot g_{\text{ml}} \cdot \mathbb{1}_{\text{irrigation active}}$ (flow rate × gain coefficient)
\end{itemize}

\subsubsection{Digital Twin Controller}

The Digital Twin maintains state variables:

\begin{equation}
\mathbf{s}(t) = [\text{moisture}(t), \text{temperature}(t), \text{humidity}(t), \text{rain\_pred}(t), \text{irr\_active}(t)]
\end{equation}

Control decisions update via the evaluate\_irrigation() method, which applies hysteresis logic to the baseline threshold model.

\subsubsection{ANFIS Model Structure}

The ANFIS system consists of:
\begin{itemize}
    \item \textbf{Fuzzification Layer}: Triangular and trapezoidal membership functions for moisture (dry, medium, wet), temperature (cold, moderate, hot), humidity (low, medium, high), and rain (no, yes).
    \item \textbf{Inference Layer}: Fuzzy rules evolved via genetic algorithm to maximize prediction accuracy on training data.
    \item \textbf{Defuzzification Layer}: Center-of-gravity method to produce scalar irrigation probability output.
\end{itemize}


\section{Implementation}
\label{sec:ch3sec3}

\subsection{Technologies and Libraries}

\subsubsection{Backend Stack}

\begin{itemize}
    \item \textbf{FastAPI 0.104+}: High-performance asynchronous web framework for REST API endpoints with automatic OpenAPI documentation.
    \item \textbf{Python 3.10+}: Primary implementation language for simulation and machine learning modules.
    \item \textbf{Pydantic}: Data validation and serialization for API request/response payloads.
    \item \textbf{multiprocessing.Pool}: Parallel execution of ANFIS fitness evaluation during genetic algorithm training.
\end{itemize}

\subsubsection{Frontend Stack}

\begin{itemize}
    \item \textbf{SAPUI5 1.120+}: Enterprise-grade UI framework from SAP providing accessibility, theming, and responsive design out-of-the-box.
    \item \textbf{sap.viz.ui5}: Visualization components for interactive line charts displaying moisture and water usage trends.
    \item \textbf{sap.ui.core.mvc}: Model-View-Controller pattern for separation of concerns between data models and UI rendering.
    \item \textbf{JSONModel}: Client-side JSON data binding for reactive UI updates without server round-trips.
\end{itemize}

\subsubsection{Machine Learning Stack}

\begin{itemize}
    \item \textbf{NumPy 1.24+}: Numerical computations and array operations.
    \item \textbf{scikit-learn 1.3+}: Preprocessing and model evaluation utilities.
    \item \textbf{Custom ANFIS Implementation}: In-house fuzzy inference engine with scikit-fuzzy-like API for research flexibility.
\end{itemize}

\subsection{Dataset}

\subsubsection{Scenario Generation}

Experimental scenarios are procedurally generated to ensure reproducibility and systematic coverage of environmental conditions:

\begin{itemize}
    \item \textbf{Steps per scenario}: 1000 time steps (representing 1000 minutes $\approx$ 16.7 hours each)
    \item \textbf{Rain probability}: 12\% per step, generating realistic precipitation patterns
    \item \textbf{Temperature range}: $U(18°C, 36°C)$ uniform distribution
    \item \textbf{Humidity range}: $U(30\%, 90\%)$ uniform distribution
    \item \textbf{Seeding}: Deterministic generation via numpy.Random(seed) for experimental reproducibility
\end{itemize}

\subsubsection{Training Data for ANFIS}

\begin{itemize}
    \item \textbf{Train set}: 500 samples of (moisture, temperature, humidity, rain\_prediction) → irrigation\_decision pairs
    \item \textbf{Test set}: 200 samples for model validation and generalization assessment
    \item \textbf{Feature scaling}: All inputs normalized to [0,1] range for fuzzy membership functions
    \item \textbf{Target labels}: Binary irrigation decisions (0: no irrigation, 1: irrigate) derived from baseline controller on diverse scenarios
\end{itemize}

\subsection{Model Architecture}

\subsubsection{ANFIS Genetic Algorithm Training}

\begin{algorithm}
\caption{ANFIS-GA Training}
\label{alg:anfis-ga}
\begin{algorithmic}[1]
\REQUIRE Dataset $\mathcal{D}_{\text{train}}$, population size $P = 40$, generations $G = 60$
\ENSURE Evolved ANFIS model with optimized fuzzy rules

\STATE Initialize population of $P$ random rule sets
\FOR{$g = 1$ to $G$}
    \FOR{each individual in population}
        \STATE Evaluate fitness: $f = \text{accuracy}(\text{individual}, \mathcal{D}_{\text{train}})$
    \ENDFOR
    \STATE Select top 50\% by fitness (tournament selection)
    \STATE Apply crossover (single-point) with probability 0.8
    \STATE Apply mutation (bit-flip on rule outputs) with probability 0.1
    \STATE Elitism: preserve top 2 individuals
\ENDFOR
\RETURN Best individual (evolved rule set)
\end{algorithmic}
\end{algorithm}

\subsubsection{Parallel Fitness Evaluation}

To accelerate training, fitness evaluation is parallelized across CPU cores:

\begin{equation}
\text{fitness}_{\text{parallel}} = \text{Pool.map}(\text{evaluate\_individual}, \text{population})
\end{equation}

This reduces training time from $O(G \cdot P)$ sequential evaluations to $O(\lceil G \cdot P / N_{\text{cores}} \rceil)$.

\subsection{Metrics}

The system computes the following performance metrics:

\subsubsection{Irrigation Efficiency}

\begin{equation}
\text{irrigation\_efficiency} = \frac{\text{irrigation events (baseline)}}{\text{irrigation events (candidate)}}
\end{equation}

\subsubsection{Water Usage}

\begin{equation}
\text{total water usage (L)} = \text{irrigation steps} \times \text{flow rate (mL/min)} / 1000
\end{equation}

\subsubsection{Sampling Accuracy}

\begin{equation}
\text{accuracy} = \frac{\text{matching decisions}}{\text{total steps}} \times 100\%
\end{equation}

where matching decisions are timesteps where sparse and continuous sampling agree on irrigation state.

\subsubsection{ANFIS Generalization Metrics}

\begin{itemize}
    \item Mean Squared Error: $\text{MSE} = \frac{1}{N} \sum_{i=1}^{N} (\hat{p}_i - p_i)^2$
    \item Classification Accuracy: $\text{Acc} = \frac{\text{correct predictions}}{N} \times 100\%$
    \item Predicted Probability Statistics: mean, min, max across test set
\end{itemize}

\subsection{Statistical Analysis}

\subsubsection{Experimental Design}

Each experiment is executed with:
\begin{itemize}
    \item Fixed scenario seed for determinism
    \item Parameter variations: threshold $\in [10, 80]$, hysteresis $\in [0, 20]$, flow rate $\in [1, 10000]$ mL/min
    \item Result caching to prevent redundant computations
\end{itemize}

\subsubsection{Hypothesis Testing}

\begin{itemize}
    \item \textbf{H1}: Sparse sampling (interval $\geq 10$) maintains $\geq 95\%$ decision accuracy vs. continuous sampling
    \item \textbf{H2}: ANFIS-GA irrigation produces statistically equivalent or superior water use efficiency compared to threshold baseline
    \item \textbf{H3}: ANFIS learned policies generalize to unseen environmental scenarios with $\geq 80\%$ test accuracy
\end{itemize}

\subsubsection{Performance Comparisons}

Statistical significance is assessed via:
\begin{itemize}
    \item Paired t-tests for mean water usage differences
    \item Chi-squared tests for decision agreement rates
    \item Confidence intervals (95\%) on accuracy metrics
\end{itemize}

\section{API and Integration}

\subsection{REST API Endpoints}

\begin{itemize}
    \item \texttt{GET /api/experiment}: Execute baseline threshold experiment
    \item \texttt{GET /api/experiment/sampling}: Execute sparse sampling comparison
    \item \texttt{GET /api/experiment/anfis}: Execute ANFIS-GA intelligence experiment
    \item Query parameters: \texttt{steps}, \texttt{threshold}, \texttt{hysteresis}, \texttt{flow\_rate\_ml}, \texttt{seed}
\end{itemize}

\subsection{Frontend Integration}

The SAPUI5 frontend consumes these endpoints asynchronously, caching results client-side to minimize redundant API calls and provide rapid parameter exploration. Charts dynamically bind to experiment entries, displaying moisture trends and water usage in mL per time step.

